\chapter{Metriken auf Funktionenräumen}
\label{b47guide:ch:function-spaces}

Die folgenden Aussagen verwenden auf Funktionenmengen ausschließlich
metrische Strukturen. Addition und Skalarmultiplikation von Funktionen
werden nicht vorausgesetzt.

\section{Gleichmäßige Konvergenz und Vollständigkeit}

\begin{definition}[Gleichmäßige Konvergenz]
\label{b47guide:def:uniform-convergence}
Seien $X$ eine Menge, $(Y,d_Y)$ ein metrischer Raum und
$f_n,f:X\to Y$. Die Folge $(f_n)$ heißt \emph{gleichmäßig gegen $f$
konvergent}, wenn
\[
 \forall\varepsilon>0\ \exists N\ \forall n\geq N\ \forall x\in X:
 \quad d_Y(f_n(x),f(x))<\varepsilon.
\]
Bei \emph{punktweiser Konvergenz} darf $N$ zusätzlich von $x$ abhängen.
\end{definition}

\begin{theorem}[Gleichmäßiges Cauchy-Kriterium]
\label{b47guide:thm:uniform-cauchy}
Seien $S$ eine Menge, $(Y,d_Y)$ ein vollständiger metrischer Raum und
$f_n:S\to Y$ Abbildungen. Es gibt genau dann eine Abbildung $f:S\to Y$
mit $f_n\to f$ gleichmäßig, wenn
\[
 \forall\varepsilon>0\ \exists N\ \forall m,n\geq N\ \forall s\in S:
 \quad d_Y(f_m(s),f_n(s))<\varepsilon.
\]
Der Grenzwert ist eindeutig. Beschränktheit der Abbildungen wird nicht
verlangt. Bei bereits gegebenem punktweisen Grenzwert genügt die
angegebene Bedingung auch ohne Vollständigkeit von $Y$ für gleichmäßige
Konvergenz. \cite{uniform-cauchy}
\end{theorem}
\begin{theorem}[Stetigkeit gleichmäßiger Grenzwerte]
\label{b47guide:thm:uniform-limits}
Seien $X,Y$ metrische Räume und sei $f_n:X\to Y$ eine gleichmäßig
gegen $f:X\to Y$ konvergente Folge.
Sind alle $f_n$ stetig, so ist $f$ stetig.
Sind alle $f_n$ gleichmäßig stetig, so ist $f$ gleichmäßig stetig.
Sind alle $f_n$ $L$-Lipschitz mit derselben Konstante $L$, so ist
$f$ $L$-Lipschitz; für diese letzte Aussage genügt bereits
punktweise Konvergenz.
\end{theorem}

\begin{definition}[Supremumsmetrik auf beschränkten Abbildungen]
\label{b47guide:def:sup-metric}
Sei $X$ eine Menge und sei $(Y,d_Y)$ ein nichtleerer metrischer Raum.
Es bezeichne $B(X,Y)$ die Menge aller Abbildungen $f:X\to Y$ mit
beschränktem Bild. Für $f,g\in B(X,Y)$ setzen wir
\[
 d_\infty(f,g)=\sup_{x\in X}d_Y(f(x),g(x)).
\]
Für $X=\varnothing$ wird $d_\infty(f,g)=0$ gesetzt.
Ist $X$ selbst metrisch, bezeichnet $C_b(X,Y)$ die Teilmenge der
stetigen beschränkten Abbildungen.
Das Supremum dient hier ausschließlich zur Definition eines Abstands.
\end{definition}

\begin{theorem}[Vollständigkeit beschränkter metrischer Funktionenräume]
\label{b47guide:thm:function-space-completeness}
Die Vorschrift $d_\infty$ ist eine endliche Metrik auf $B(X,Y)$.
Konvergenz bezüglich dieser Metrik ist genau gleichmäßige Konvergenz
innerhalb von $B(X,Y)$.
Ist $Y$ vollständig, so ist $(B(X,Y),d_\infty)$ vollständig.
Ist $X$ metrisch, so ist $C_b(X,Y)$ abgeschlossen in $B(X,Y)$
und bei vollständigem $Y$ ebenfalls vollständig.
\cite{metric-notes}
\end{theorem}

\section{Kompaktheit von Familien stetiger Abbildungen}

\begin{definition}[Gleichgradige Stetigkeit und punktweise relative Kompaktheit]
\label{b47guide:def:equicontinuity}
Seien $K,Y$ metrische Räume und sei $\mathcal F$ eine Familie
stetiger Abbildungen $K\to Y$.
Die Familie heißt \emph{gleichgradig stetig}, wenn für jedes $x\in K$
und jedes $\varepsilon>0$ ein $\delta>0$ existiert mit
\[
 d_K(x,z)<\delta\quad\Longrightarrow\quad
 d_Y(f(x),f(z))<\varepsilon
 \qquad(z\in K,\ f\in\mathcal F).
\]
Sie heißt \emph{gleichmäßig gleichgradig stetig}, wenn $\delta$
auch unabhängig von $x$ gewählt werden kann.
Sie heißt \emph{punktweise relativ kompakt}, wenn für jedes $x\in K$
der Abschluss von $\{f(x):f\in\mathcal F\}$ in $Y$ kompakt ist.
\end{definition}

\begin{theorem}[Gleichgradige Stetigkeit auf kompakten Räumen]
\label{b47guide:thm:uniform-equicontinuity}
Ist $K$ kompakt, so ist jede gleichgradig stetige Familie von
Abbildungen $K\to Y$ gleichmäßig gleichgradig stetig.
\end{theorem}

\begin{theorem}[Arzelà--Ascoli in metrischer Form]
\label{b47guide:thm:ascoli}
Sei $K$ ein kompakter metrischer Raum und sei $Y$ ein nichtleerer
vollständiger metrischer Raum. Es bezeichne $C(K,Y)$ die Menge
aller stetigen Abbildungen $K\to Y$, versehen mit der
Supremumsmetrik $d_\infty$. Alle diese Abbildungen sind beschränkt,
und $C(K,Y)$ ist vollständig.
Für $\mathcal F\subseteq C(K,Y)$ sind äquivalent:
\begin{enumerate}
\item Der Abschluss von $\mathcal F$ in $C(K,Y)$ ist kompakt.
\item $\mathcal F$ ist bezüglich $d_\infty$ total beschränkt.
\item $\mathcal F$ ist gleichgradig stetig und punktweise relativ kompakt.
\end{enumerate}
In der dritten Bedingung kann punktweise relative Kompaktheit
durch totale Beschränktheit jeder Wertemenge
$\{f(x):f\in\mathcal F\}$ in $Y$ ersetzt werden.
\cite{ascoli-notes}
\end{theorem}

\begin{theorem}[Teilfolgenfassung und Kompaktheitskriterium]
\label{b47guide:thm:ascoli-sequences}
Unter den Voraussetzungen von Satz~\ref{b47guide:thm:ascoli} besitzt
jede Folge aus einer gleichgradig stetigen und punktweise relativ
kompakten Familie eine gleichmäßig konvergente Teilfolge; ihr
Grenzwert gehört zu $C(K,Y)$.
Eine Familie $\mathcal F\subseteq C(K,Y)$ ist genau dann kompakt,
wenn sie in $C(K,Y)$ abgeschlossen, gleichgradig stetig und
punktweise relativ kompakt ist.
\end{theorem}
\section{Darstellung durch Abstandsfunktionen}

\begin{theorem}[Kuratowski--Fréchet-Einbettung]
\label{b47guide:thm:kuratowski-frechet}
Sei $(X,d)$ ein nichtleerer metrischer Raum und sei $x_0\in X$ fest.
Für $x\in X$ definiert
\[
 J(x)(z)=d(x,z)-d(x_0,z)\qquad(z\in X)
\]
eine beschränkte reellwertige Funktion, und es gilt
\[
 |J(x)(z)|\leq d(x,x_0),\qquad
 d_\infty(J(x),J(y))=d(x,y)\qquad(x,y,z\in X).
\]
Damit ist $J:X\to B(X,\mathbb R)$ eine isometrische Einbettung in einen
vollständigen metrischen Raum. Der Abschluss von $J[X]$ in diesem Raum
ist eine Vervollständigung von $X$.
Ist $D\subseteq X$ dicht, so gelten dieselben Aussagen für die
Einschränkungen $J(x)|_D$ in $B(D,\mathbb R)$. Bei separablem $X$ kann
$D$ höchstens abzählbar gewählt werden. \cite{kuratowski-notes}
\end{theorem}